%----------------------------------------------------------------------------------------
%	PACKAGES AND DOCUMENT CONFIGURATIONS
%----------------------------------------------------------------------------------------

% Original copied from https://www.overleaf.com/latex/templates/template-for-foundation-master-project-reports/dtzztzckwswh

\documentclass[12pt]{article}
% Times New Roman
\usepackage{tgtermes}

% Adjusting margins to personal my need
\addtolength{\oddsidemargin}{-.5in}
\addtolength{\evensidemargin}{-.5in}
\addtolength{\textwidth}{1in}
\addtolength{\topmargin}{-.5in}
\addtolength{\textheight}{1in}

% Graphics
\usepackage{graphicx}
\graphicspath{{figures/}}

% Math
\usepackage{amssymb}
\usepackage{amsmath} % Required for some math elements 

% Other
\usepackage{algorithmic}
\usepackage{array}
\usepackage[caption=false,font=footnotesize]{subfig}
\usepackage{lipsum}
\usepackage{hyperref}




%----------------------------------------------------------------------------------------
%	MAIN PART
%----------------------------------------------------------------------------------------
\begin{document}

\title{RPMI coursework title} % Title
%\author{Konstantin Akhmadeev, That Guy, Yet Another Guy (in alphabetic order)}
\author{Daniel Sikar}
\date{\today} % Date for the report
\maketitle % Inserts the title, author and date

% TODO
% 1. Sections

\begin{abstract}
%% Text of the abstract
This template will help you to write your bibliographic and final reports using \LaTeX{}. You'll find here the examples of text structuring as well as tables, figures, citations and references. For other features of \LaTeX, see tutorials on \href{https://www.overleaf.com/learn}{\textbf{Overleaf}} or use this \href{https://wch.github.io/latexsheet/}{\textbf{cheatsheet}}. To work with this template, download its entire folder (including /bibliography and /figures), and run your \LaTeX{}  editor like \href{http://www.xm1math.net/texmaker/}{\textbf{Texmaker}} or \href{https://www.overleaf.com}{\textbf{Overleaf}}. Then make a plan by changing the document's structure with \textit{section} and \textit{subsection} commands. Finally, delete the \textit{lipsum} fillings and start writing you report. Good luck!
\end{abstract}


\section{Introduction}

\textbf{Example of citation: \cite{Smith_2013} and \cite{Smith_2012}}. \lipsum[2]

\section{Main part}
\label{S:1}
\lipsum[9]

Examples of lists:
\begin{itemize}
\item Bullet point one
\item Bullet point two
\end{itemize}

\begin{enumerate}
\item Numbered list item one
\item Numbered list item two
\end{enumerate}

\subsection{Subsection 1}

\textbf{Example of table reference: see Table \ref{tab:example}}.
\lipsum[4]

\begin{table}[h] 
\centering
\begin{tabular}{l l l}
\hline
\textbf{Treatments} & \textbf{Response 1} & \textbf{Response 2}\\
\hline
Treatment 1 & 0.0003262 & 0.562 \\
Treatment 2 & 0.0015681 & 0.910 \\
Treatment 3 & 0.0009271 & 0.296 \\
\hline
\end{tabular}
\caption{Table caption}
\label{tab:example}
\end{table}

\subsection{Subsection 2}

\textbf{Example of figure reference: see Figure \ref{fig:example}}. 
\lipsum[5]

\begin{figure}[ht]
\centering\includegraphics[width=0.4\linewidth]{placeholder}
\caption{Figure caption}
\label{fig:example}
\end{figure}

\textbf{Example of equation reference: see Equation \eqref{eq:emc}}. 
\lipsum[6]

\begin{equation} 
\label{eq:emc}
e = mc^2
\end{equation}

\subsection{Subsection 3}

\textbf{Example of reference to Section \ref{S:1}.} 
\lipsum[7]
\lipsum[8]



\section{Results (optional)}

\lipsum[9]

\section{Conclusion}

\lipsum[10]

%----------------------------------------------------------------------------------------
%	Bibliography
%----------------------------------------------------------------------------------------
\bibliography{bibliography/sample}{}
\bibliographystyle{plain}

\end{document}